% -----------------------------------------------
% MIREX 2026 -- Music Performance Difficulty Prediction
% CrescendAI technical report (ISMIR 2026 LBD format)
% -----------------------------------------------
\documentclass{article}
\usepackage[T1]{fontenc}
\usepackage[utf8]{inputenc}
\usepackage[lbd]{ismir}
\usepackage{amsmath,cite,url}
\usepackage{graphicx}
\usepackage{color}
\usepackage{booktabs}
\def\lbd{}

\title{Transcription-First Piano Difficulty Estimation with a Fine-Tuned
Symbolic Foundation Model}

\oneauthor
  {Jai Dhiman}
  {CrescendAI \\ {\tt jaidhiman2000@gmail.com}}

\def\authorname{J. Dhiman}

\sloppy

\begin{document}

\maketitle

\begin{abstract}
We describe the CrescendAI submission to the MIREX 2026 Music Performance
Difficulty Prediction task. The system is transcription-first: audio is
converted to MIDI with Transkun, the resulting symbolic sequence is encoded by
MoonBeam-839M fine-tuned with LoRA adapters, token embeddings are mean-pooled
over the whole piece, and a ridge head emits one real-valued difficulty score.
The submitted system averages the predictions of five adapters that differ only
by training seed. Training used 7{,}665 transcriptions of the public Piano
Syllabus corpus; no other data was used, and the evaluation repertoire is
disjoint from that corpus by the task's own construction. On composer-disjoint
cross-validation the recipe scores Kendall's $\tau_c$ 0.8430 for a single
adapter and 0.8461 for the five-adapter ensemble, against 0.8068 for a
supervision-matched hand-feature baseline on identical folds. We also estimate
the ceiling imposed by label noise at $\tau_c$ 0.8668, leaving roughly $+0.024$
of headroom on this corpus. We report no score for the submitted model itself:
it is trained on every labelled piece available to us, so any evaluation of it
would be train-on-test.
\end{abstract}

\section{System Overview}\label{sec:overview}

The container implements \texttt{WAV path} $\rightarrow$ \texttt{float}. Four
stages run per item:

\begin{enumerate}
\itemsep0em
\item \textbf{Transcription.} Transkun~\cite{transkun} converts the recording
      to a MIDI note sequence with onsets, offsets, velocities, and sustain
      pedal.
\item \textbf{Encoding.} The sequence is tokenised and passed through
      MoonBeam-839M~\cite{moonbeam2025}, a MIDI-native causal language model,
      carrying LoRA adapters~\cite{hu2022lora}. Sequences longer than the
      1024-token context are split into consecutive chunks and every chunk is
      encoded; encoder cost is therefore linear in piece duration.
\item \textbf{Pooling.} Token embeddings are averaged over the whole piece.
\item \textbf{Regression.} A ridge head maps the pooled embedding to a scalar.
\end{enumerate}

We chose transcription-first over an audio-native model after a frozen
bake-off in which MoonBeam-839M with mean pooling reached $\tau_c$ 0.8257
against 0.7790 for a competing symbolic encoder on identical folds. A control
established that the gain is backbone-specific rather than an artefact of the
pooling operator.

\section{Training Data}\label{sec:data}

\textbf{Dataset.} Piano Syllabus
(PSyllabus)~\cite{ramoneda2025audio,psyllabus2025dataset}, Zenodo
\texttt{10.5281/zenodo.14794592}: 7{,}901 recordings of solo piano repertoire,
graded on an 11-level scale aggregated from 18 grading systems. No other corpus
contributed to training.

\textbf{Size actually used.} PSyllabus distributes labels and video identifiers
rather than audio, so recordings were fetched individually and transcribed
locally. We obtained and transcribed \textbf{7{,}665 of 7{,}901 pieces
(97.0\%)}; the remaining 236 are no longer retrievable (186 removed, 48
private, 2 region-blocked). The five submitted adapters were each trained on
all 7{,}665, with nothing held out.

\textbf{Preprocessing.} Audio is transcribed to MIDI with Transkun; no manual
cleaning, filtering, or repertoire curation was applied. Training samples a
single random 1024-token window per piece per epoch, while inference encodes
every chunk and mean-pools, so the training and inference views deliberately
differ. We swept the training window length and retained 1024.

\textbf{Contamination control.} The task specifies that the evaluation
repertoire is drawn from sources outside existing public piano difficulty
datasets. PSyllabus is such a dataset, so disjointness holds by construction.
We disclose the corpus in full so that organisers can verify this directly; at
the time of writing the announced list of forbidden pieces and composers had
not been published, so no intersection against it could be computed. An
automated check is held ready to run against that list should it appear.

\section{Model and Label Space}\label{sec:model}

\textbf{Architecture.} MoonBeam-839M with LoRA on the top 5 of 15 transformer
layers, adapting the seven projection matrices in each
($\{q,k,v,o,\mathrm{gate},\mathrm{up},\mathrm{down}\}$, 35 modules). The
backbone is otherwise frozen. Training ran 3 epochs at learning rate
$10^{-4}$.

\textbf{Label space.} PSyllabus grades are ordinal with 11 levels, and the
published grade is the rounded mean of the judgements of one or more grading
boards ($r = 0.981$ against the raw board mean). We apply \emph{no} remapping:
the head is fit by ridge regression directly on the integer grade and emits an
unrounded real value, because the task's metric is rank-based and its own scale
is undisclosed. Emitting a continuous score avoids committing to a granularity
we cannot observe.

\textbf{Post-processing.} The submitted score is the unweighted mean of the
five member predictions. Members differ only by training seed (2026--2030),
which varies both LoRA initialisation and window sampling. Scores are neither
clipped nor rounded. On an item that fails to transcribe or encode, the
container emits the corpus median and logs the failure to \texttt{stderr}; this
bounds the effect of unreadable inputs on the run as a whole rather than
aborting it.

\section{Validation}\label{sec:validation}

All figures below come from composer-disjoint 5-fold cross-validation over 900
evaluation pieces spanning 900 distinct composers, with byte-identical folds
across arms and paired bootstrap over pieces (2000 resamples, seed 2026).

\begin{table}[h]
\centering
\begin{tabular}{lc}
\toprule
arm (identical folds, $n=900$) & $\tau_c$ \\
\midrule
37 hand features, ridge (matched) & 0.8068 \\
MoonBeam-839M frozen, mean-pool & 0.8257 \\
MoonBeam-839M + LoRA, single adapter & 0.8430 \\
\textbf{5-adapter prediction mean (recipe)} & \textbf{0.8461} \\
\midrule
\emph{estimated label ceiling} & \emph{0.8668} \\
\bottomrule
\end{tabular}
\end{table}

The fine-tuned encoder beats the supervision-matched hand-feature baseline by a
margin whose 95\% bootstrap interval excludes zero, and the five-adapter
ensemble beats a single adapter by $+0.0051$ CI95$[+0.0019,+0.0082]$. The
per-adapter spread (std 0.0033) is the same order as the ensemble gain, which
is consistent with variance averaging rather than any member holding unique
signal.

\textbf{Label ceiling.} Fitting a graded response model~\cite{samejima1969grm}
to all 13{,}659 raw board judgements, with per-board discriminations and
thresholds, gives the $\tau_c$ that a \emph{perfect} model would score against
this label: \textbf{0.8668} CI95$[0.8549, 0.8805]$. Total remaining headroom on
this corpus is therefore about $+0.024$. A posterior-predictive check indicates
the fit understates board agreement, so the estimate is conservative. This
ceiling is specific to this corpus's board composition; the MIREX test set is
independently curated and its own ceiling is unmeasurable from here.

\textbf{What we do not claim.} Published results on PSyllabus are not
comparable to the figures above: they use random rather than composer-disjoint
splits, a different subset, and a different grade mix. We make no claim of
improving on them. We likewise report no $\tau_c$ for the submitted model,
because every labelled piece is in its training set.

\section{Compute}\label{sec:compute}

Measured on one NVIDIA L4 over 8 recordings spanning 15--1806\,s, 0/8 failures:
the five-member ensemble costs $5.00\,\mathrm{s} + 0.0410 \times
\mathrm{duration}$ per item, against $4.93\,\mathrm{s} + 0.0360 \times
\mathrm{duration}$ for a single adapter. An earlier run of the same protocol on
a separately built copy of the model gave $4.97\,\mathrm{s} + 0.0400 \times
\mathrm{duration}$, so the figure reproduces. At our corpus's median piece length of
103\,s this is an \textbf{average inference time of 9.2\,s per piece}, or
roughly 9{,}400 items per 24 hours on one GPU; at the mean duration of our
evaluation set (150\,s) the average is 11.2\,s per piece. The ensemble's cost
falls almost entirely on the slope rather than the fixed term, because the
encoder processes every chunk of a piece.

\textbf{Peak GPU memory is 5{,}976\,MiB} (roughly 5.8\,GiB), sampled
device-wide at 0.5\,s intervals across the whole run. We report the device
figure rather than the framework's own counters because the pipeline spans two
processes: the encoder's allocator reserved 3{,}858\,MiB, and the remainder is
CUDA context plus the Transkun transcriber, which runs in a separate process
with its own context. A framework-only number would have omitted the
transcriber entirely. An 8\,GiB card is sufficient.

The submitted image is self-contained and performs no network access at
inference time.

\bibliography{references}

\end{document}
